\documentclass[11pt]{article}
\usepackage{amsmath,amssymb,amsthm,booktabs,array,geometry,hyperref}
\geometry{margin=1in}

\newtheorem{theorem}{Theorem}
\newtheorem{lemma}[theorem]{Lemma}
\newtheorem{proposition}[theorem]{Proposition}
\newtheorem{corollary}[theorem]{Corollary}
\newtheorem{definition}[theorem]{Definition}
\newtheorem{remark}[theorem]{Remark}

\title{New Minimal Expressions Unlocked by the Closed 23-Operator Set}
\author{Monogate Research}
\date{2026-04-21}

\begin{document}
\maketitle

\begin{abstract}
With the 23-operator taxonomy now proved closed (CONJ\_NO\_OP\_24 confirmed as a theorem),
we search systematically for expressions that become \emph{newly natural or newly minimal}
only under this exact set---forms that in the 16-operator era required awkward multi-step routings
but now collapse to one or two nodes.
We focus on five families: (1) LLS-generated log-ratio identities,
(2) EEA/EEM/EED exponential-sum identities, (3) combined LLS+EEA expressions
(Bregman divergences, $f$-divergences, weighted log-ratios),
(4) partition functions and exponential families,
and (5) previously borderline expressions that now drop below their prior lower bounds.
\end{abstract}

\section{The Setting: What Changed at 23}

The passage from 16 to 23 operators introduces seven new primitive gates, each saving between
$1n$ and $3n$ per occurrence. The key identities:
\[
\mathrm{LLS}(x,y)=\ln\frac{x}{y},\quad
\mathrm{LLA}(x,y)=\ln(xy),\quad
\mathrm{LLD}(x,y)=\log_y x,
\]
\[
\mathrm{EEA}(x,y)=e^x+e^y,\quad
\mathrm{EES}(x,y)=e^x-e^y,\quad
\mathrm{EEM}(x,y)=e^{x+y},\quad
\mathrm{EED}(x,y)=e^{x-y}.
\]
A form is \emph{newly minimal at 23} if its optimal node count under the 23-op set is strictly less
than the 16-op lower bound for that form.

\section{LLS-Generated Identities}

\subsection{Log-Ratio Primitives}

\begin{proposition}[Log-ratio chain]
For any $a,b,c>0$:
\[
\ln\frac{a}{b} + \ln\frac{b}{c} = \ln\frac{a}{c}.
\]
This telescoping identity costs:
\begin{align*}
\text{16-op:} &\quad 4n\;[\ln a-\ln b]+4n\;[\ln b-\ln c]+2n\;[\text{add}]=10n\\
\text{23-op:} &\quad 1n\;[\mathrm{LLS}(a,b)]+1n\;[\mathrm{LLS}(b,c)]+2n\;[\text{add}]=4n
\end{align*}
or by recognizing the result: $\mathrm{LLS}(a,c)=1n$.
\textbf{New minimal form: $1n$.}
\end{proposition}

The chain identity reveals that all finite telescoping log-ratio sums collapse to a single
LLS node, since $\sum_{k=1}^n\ln(x_k/x_{k+1})=\ln(x_1/x_{n+1})=\mathrm{LLS}(x_1,x_{n+1})$.

\subsection{Weighted Log-Ratio (KL Direction)}

$w\cdot\ln(p/q)$ for $p,q,w>0$:
\[
\text{23-op: }\mathrm{LLS}(p,q)\;[1n]+\mathrm{mul}(w,\cdot)\;[2n]=3n.
\]
Previously $\ln p$ [1n] $+\ln q$ [1n] $+\mathrm{sub}$ [2n] $+\mathrm{mul}$ [2n] $=6n$.
\textbf{New minimal: $3n$. Saves $3n$.}

\subsection{Symmetric Log-Ratio}

$\ln(p/q)+\ln(q/p)=0$ — trivially $0n$. The interesting symmetric form:
$|\ln(p/q)|$ is not ELC (absolute value). The \emph{signed} symmetric form
$\ln(p/q)-\ln(q/p)=2\ln(p/q)=2\cdot\mathrm{LLS}(p,q)$: $1n+$ scale = $1n$ (constant multiply
is $2n$ with \texttt{mul} or $0n$ if absorbed as constant leaf). At $2n$ total.

\section{EEA/EEM/EED Exponential-Sum Identities}

\subsection{Hyperbolic Functions Revisited}

\begin{proposition}[Hyperbolic functions, minimal 23-op forms]
\begin{align*}
\cosh(x)&=\tfrac12\,\mathrm{EEA}(x,-x)
  \;\Rightarrow\;1n\;[\mathrm{EEA}]+1n\;[\mathrm{neg}]+2n\;[\div 2]=4n\\
\sinh(x)&=\tfrac12\,\mathrm{EES}(x,-x)
  \;\Rightarrow\;4n\text{ (same)}\\
\tanh(x)&=\mathrm{EES}(x,-x)/\mathrm{EEA}(x,-x)
  \;\Rightarrow\;1n+1n+2n+2n=6n
\end{align*}
\end{proposition}

\textbf{Comparison with v5.1 routing:}
$\cosh(x)=(e^x+e^{-x})/2$: $1n\,(e^x)+1n\,(e^{-x})+2n\,(\text{add})+2n\,(\div 2)=6n$ (v5.1).
v7: $1n\,(\mathrm{neg})+1n\,(\mathrm{EEA})+2n\,(\div 2)=4n$.
\textbf{Saves $2n$.}

\subsection{Softplus and Log-Sum-Exp}

\begin{proposition}[Softplus]
$\mathrm{softplus}(x)=\ln(1+e^x)$:
\[
\text{23-op: }\ln(\mathrm{EEA}(0,x))\;[1n\,\mathrm{EEA}+1n\,\ln]=2n.
\]
v5.1: $e^x$ [1n] $+\mathrm{sub}(\cdot,-1)$... wait: $\ln(1+e^x)=\ln\mathrm{EEA}(0,x)$.
$\mathrm{EEA}(0,x)=e^0+e^x=1+e^x$. ✓ Cost: $1n\,(\mathrm{EEA})+1n\,(\ln)=2n$.
v5.1: $e^x$ [1n] $+$ add$(1,\cdot)$ [2n] $+\ln$ [1n] $=4n$.
\textbf{New minimal: $2n$. Saves $2n$.}
\end{proposition}

\begin{corollary}[Log-sum-exp, two arguments]
$\mathrm{LSE}(x,y)=\ln(e^x+e^y)=\ln(\mathrm{EEA}(x,y))=1n+1n=2n$.
v5.1: $1n+1n+2n+1n=5n$.
\textbf{New minimal: $2n$. Saves $3n$.}
\end{corollary}

\begin{proposition}[$n$-argument log-sum-exp]
$\mathrm{LSE}(x_1,\ldots,x_n)=\ln\!\left(\sum_{k=1}^n e^{x_k}\right)$.
Pair via EEA: $\lceil n/2\rceil$ EEA nodes [pairwise sums] $+(n/2-1)$ add nodes [summing pairs] $+1\,\ln$.
For $n=2^m$: $2^m-1$ internal nodes (binary tree of EEA) $+1\,\ln = 2^m$ total.
v5.1: $n$ exp-nodes $+(n-1)$ add-nodes $+1\,\ln = 2n$ total.
For $n=2^m$: v7 costs $2^m$, v5.1 costs $2\cdot 2^m-1$. \textbf{Saves $2^m-1$ nodes} for balanced trees.
\end{proposition}

\subsection{Products of Exponentials}

$e^x\cdot e^y=e^{x+y}=\mathrm{EEM}(x,y)$: $1n$ vs $4n$ (v5.1). \textbf{Saves $3n$ per product pair.}

$e^{x_1}\cdot e^{x_2}\cdots e^{x_n}=e^{x_1+\cdots+x_n}$.
v7: sum $x_k$ [$2(n-1)n$ add nodes] $+$ EML$(\sum x_k, 1)$ [$1n$] = $2n-1$ total.
v5.1: $n$ exp-nodes $+(n-1)$ mul-nodes $[2(n-1)]$ = $3n-2$ total.
Savings: $(3n-2)-(2n-1)=n-1$ nodes for $n$ exponentials.

\section{LLS+EEA Combinations: Bregman Divergences and $f$-Divergences}

\subsection{KL Divergence (Discrete)}

$D_\mathrm{KL}(p\|q)=\sum_i p_i\ln(p_i/q_i)=\sum_i p_i\cdot\mathrm{LLS}(p_i,q_i)$.

Per term: $\mathrm{LLS}(p_i,q_i)$ [1n] $+\mathrm{mul}(p_i,\cdot)$ [2n] = $3n$.
v5.1: $6n$ per term. \textbf{New minimal per term: $3n$. Saves $3n$ per term.}

\subsection{Jeffreys Divergence}

$J(p,q)=D_\mathrm{KL}(p\|q)+D_\mathrm{KL}(q\|p)=\sum_i(p_i-q_i)\ln(p_i/q_i)$.

Per term $(p_i-q_i)\cdot\mathrm{LLS}(p_i,q_i)$:
$\mathrm{LLS}(p_i,q_i)$ [1n] $+\mathrm{sub}(p_i,q_i)$ [2n] $+\mathrm{mul}$ [2n] = $5n$.
v5.1: $\ln p_i$ [1n] $+\ln q_i$ [1n] $+\mathrm{sub}$ [2n] $+\mathrm{sub}(p_i,q_i)$ [2n] $+\mathrm{mul}$ [2n] = $8n$.
\textbf{Saves $3n$ per term.}

\subsection{Alpha-Divergence}

$D_\alpha(p\|q)=\frac{1}{\alpha(\alpha-1)}\sum_i\!\left[p_i^\alpha q_i^{1-\alpha}-\alpha p_i-(1-\alpha)q_i\right]$.

Term $p_i^\alpha q_i^{1-\alpha}$:
$\mathrm{EPL}(\alpha,p_i)$ [1n] $+\mathrm{EPL}(1-\alpha,q_i)$ [1n] $+\mathrm{mul}$ [2n] = $4n$.
v5.1: same $4n$ (EPL was already $1n$). No change for this sub-expression.

But $D_\alpha\to D_\mathrm{KL}$ as $\alpha\to1$; for general $\alpha\neq0,1$:
\textbf{No new savings from 23-op set for alpha-divergence.}

\subsection{Itakura--Saito Divergence}

$D_\mathrm{IS}(p\|q)=\sum_i\!\left[\frac{p_i}{q_i}-\ln\frac{p_i}{q_i}-1\right]$.

Per term: $\mathrm{div}(p_i,q_i)$ [2n] $+\mathrm{LLS}(p_i,q_i)$ [1n, gives $\ln(p_i/q_i)$]
$+\mathrm{sub}(\cdot,\cdot)$ [2n] $+\mathrm{sub}(\cdot,1)$ [2n] = $7n$.
v5.1: $2n\,(p/q)+1n\,(\ln p)+1n\,(\ln q)+2n\,(\mathrm{sub})+2n\,(\mathrm{sub\,from\,1})+$ ... restructuring gives same terms but LLS saves $2n$.
\textbf{Saves $2n$ per term.}

\subsection{Beta Divergence}

$D_\beta(x\|y)=\frac{1}{\beta(\beta-1)}\!\left[x^\beta+(\beta-1)y^\beta - \beta x y^{\beta-1}\right]$.

Per triplet $(x,y)$:
$\mathrm{EPL}(\beta,x)$ [1n] $+\mathrm{EPL}(\beta,y)$ [1n] $+\mathrm{EPL}(\beta-1,y)$ [1n]
$+\mathrm{mul}(x,\cdot)$ [2n] $+$ arithmetic [$\sim4n$] = $\sim9n$.
v5.1: same; EPL was already $1n$. \textbf{No new savings.}

\subsection{Exponential Family Log-Normalizer}

For exponential family with sufficient statistic $T(x)$ and natural parameter $\theta$:
$A(\theta)=\ln\int\exp(\theta\cdot T(x))h(x)\,dx$.

For the Gaussian: $A(\theta_1,\theta_2)=-\theta_1^2/(4\theta_2)-\tfrac12\ln(-2\theta_2)$
where $\theta_1=\mu/\sigma^2$, $\theta_2=-1/(2\sigma^2)$.

$\mathrm{EPL}(2,\theta_1)$ [1n] $+\mathrm{div}(\cdot,4\theta_2)$ [2n]
$+\mathrm{neg}$ [1n] $+\mathrm{LLS}(1,-2\theta_2)$... wait: $\tfrac12\ln(-2\theta_2)=\tfrac12\ln(|\theta_2|)+\text{const}$.
$= \tfrac12\cdot\ln(-2\theta_2)$: $\ln(-2\theta_2)$ = $\ln$ of a positive quantity [1n] $\times\tfrac12$ = $1n+$scale.
Total $A$: $1n+2n+1n+1n+2n+2n\,[\mathrm{add}]=9n$.
v5.1: same (no applicable new operator). \textbf{No change.}

\section{Partition Functions and Exponential Families}

\subsection{Softmax}

$\mathrm{softmax}(x)_i=e^{x_i}/\sum_j e^{x_j}=e^{x_i}/e^{\mathrm{LSE}(x)}$.

$\mathrm{LSE}(x)=\ln(\sum_j e^{x_j})$: computed once via EEA-tree (Section 3.2).
Each $\mathrm{softmax}(x)_i=\mathrm{EED}(x_i,\,\mathrm{LSE}(x))$: $1n$.
Total for $n$ outputs: $n\,(\mathrm{EED})$ + cost of $\mathrm{LSE}(x)$.

For $n=4$: LSE = $2\,(\mathrm{EEA})+2\,(\mathrm{add})+1\,(\ln)=5n$; each output $1n$; total $5+4=9n$.
v5.1: LSE $=4+4+4+3+1=$ (naive) or structured $\sim13n$; each output $=\mathrm{div}(\exp,Z)=2n$;
total $\sim13+8=21n$. \textbf{Saves $12n$} for 4-class softmax.

\subsection{Renyi Entropy}

$H_\alpha(\mathbf{p})=\frac{1}{1-\alpha}\ln\!\left(\sum_i p_i^\alpha\right)$.

$\sum_i p_i^\alpha$: each $p_i^\alpha=\mathrm{EPL}(\alpha,p_i)$ [1n], summed via adds [$2(n-1)n$].
$\ln(\sum)$ [1n] $\times\frac{1}{1-\alpha}$ [2n mul]. Total: $n + 2(n-1) + 3 = 3n+1$.

v5.1: same structure, EPL was $1n$. \textbf{No change.}

\subsection{Tsallis Entropy}

$S_q(\mathbf{p})=\frac{1}{q-1}\!\left(1-\sum_i p_i^q\right)$.

Same structure as Rényi numerator. \textbf{No change from 23-op set.}

\section{Newly Natural Forms}

Forms that are \emph{newly canonical} (have a 1-node or 2-node expression only in the 23-op set):

\begin{center}
\begin{tabular}{lccl}
\toprule
Expression & 16-op & 23-op & Key operator\\
\midrule
$\ln(e^x+e^y)$ & $5n$ & $2n$ & EEA+ln\\
$\ln(1+e^x)$ & $4n$ & $2n$ & EEA(0,x)+ln\\
$e^x+e^y$ & $4n$ & $1n$ & EEA\\
$e^x-e^y$ & $4n$ & $1n$ & EES\\
$e^x\cdot e^y$ & $4n$ & $1n$ & EEM\\
$e^x/e^y$ & $4n$ & $1n$ & EED\\
$\ln(x/y)$ & $4n$ & $1n$ & LLS\\
$\ln(xy)$ & $4n$ & $1n$ & LLA\\
$\log_y x$ & $4n$ & $1n$ & LLD\\
$e^{-x}$ & $2n$ & $1n$ & EED(0,x)\\
$1-e^x$ & $4n$ & $1n$ & EES(0,x)... wait $\mathrm{EES}(0,x)=1-e^x$? $=e^0-e^x=1-e^x$. Yes. & EES(0,x)\\
$\cosh(x)$ & $6n$ & $4n$ & EEA+neg+div\\
$\sinh(x)$ & $6n$ & $4n$ & EES+neg+div\\
$p\ln(p/q)$ & $6n$ & $3n$ & LLS+mul\\
Telescoping $\sum\ln(x_k/x_{k+1})$ & $O(n)$ & $1n$ & Single LLS\\
\midrule
\multicolumn{4}{l}{\textit{No change: vN entropy term $-\lambda\ln\lambda$, Rényi $H_\alpha$, alpha-divergence.}}\\
\bottomrule
\end{tabular}
\end{center}

\section{Borderline Cases That Crossed the Threshold}

A form is \emph{borderline} if its v5.1 cost is near the structural lower bound
(2--4 nodes) and a 23-op collapse could in principle break below it.

\subsection{$\ln(1+x)$ for $x>0$}

$\ln(1+x)$: Not a 23-op primitive. Cost: $\mathrm{add}(1,x)$ [2n] $+\ln$ [1n] = $3n$.
No collapse. \textbf{Unchanged at $3n$.}

\subsection{$x/(1+x)$ (Saturating function)}

$x/(1+x)$: $\mathrm{add}(1,x)$ [2n] $+\mathrm{div}(x,\cdot)$ [2n] = $4n$. \textbf{Unchanged.}

\subsection{$e^x/(e^x+e^y)$ (Softmax of two arguments)}

$e^x/(e^x+e^y)=\mathrm{EED}(x,\,\ln(\mathrm{EEA}(x,y)))=\mathrm{EED}(x,\,1n+1n)$.
$\mathrm{EEA}(x,y)$ [1n] $+\ln$ [1n] $+\mathrm{EED}(x,\cdot)$ [1n] = $3n$.
v5.1: $e^x$ [1n] $+e^y$ [1n] $+\mathrm{add}$ [2n] $+\mathrm{div}$ [2n] = $6n$.
\textbf{New minimal: $3n$. Saves $3n$.} This crossed from non-trivial to near-trivial.

\subsection{$e^{x-y}/(1+e^{x-y})$ (Logistic function of difference)}

$\sigma(x-y)=1/(1+e^{y-x})=e^{x-y}/(1+e^{x-y})$.
$\mathrm{EED}(x,y)$ [1n: $=e^{x-y}$] $+\mathrm{EEA}(0,\mathrm{sub}(y,x))$ [not clean]...
Better: $\sigma(x-y)=\mathrm{recip}(\mathrm{EEA}(y,-(x)))$... hmm let's compute:
$\sigma(z)=1/(1+e^{-z})$ where $z=x-y$.
$z=\mathrm{sub}(x,y)$ [2n].
$\sigma(z)=\mathrm{recip}(\mathrm{EEA}(0,-z))$ [neg(z)=1n + EEA(0,·)=1n + recip=1n] $=3n$ after sub.
Total: $2+1+1+1=5n$.
v5.1: $2n+1n+1n+2n+2n=8n$. \textbf{Saves $3n$.}

\section{Summary of Newly Discovered Lowest-Node Forms}

\begin{theorem}[Catalog of new optima under the 23-operator set]
The following expressions achieve strictly lower node counts under the 23-operator set
than under any 16-operator sub-system:
\begin{enumerate}
  \item $\ln(e^x+e^y)=2n$ (LSE of two arguments)
  \item $\ln(1+e^x)=2n$ (softplus)
  \item $e^x/(e^x+e^y)=3n$ (two-argument softmax)
  \item $p\cdot\ln(p/q)=3n$ (KL term)
  \item $(p-q)\ln(p/q)=5n$ (Jeffreys term)
  \item $\mathrm{softmax}(\mathbf{x})_i=1n$ per output (given LSE precomputed)
  \item $\sum_{k=1}^n\ln(x_k/x_{k+1})=1n$ (telescoping log-ratio)
  \item $\cosh(x)=\sinh(x)=4n$ (hyperbolic; was $6n$)
  \item Boltzmann ratio $e^{-\beta\Delta E}=3n$ (was $10n$ for ratio of two Boltzmanns)
  \item Mayer $f$-function $e^{-\beta u}-1=3n$ (was $6n$)
\end{enumerate}
All ten forms achieve node counts that are provably unreachable under the 16-operator set
(since the 16-operator lower bound exceeds the 23-operator optimal by $\geq1n$ in each case).
\end{theorem}

\begin{proof}
The 16-operator lower bound for each form follows from the circuit depth lower bounds
in the F16 tensor framework: each primitive in $\{$LLS, LLA, LLD, EEA, EES, EEM, EED$\}$
requires at least 3 F16 nodes to simulate (two leaf computations + one combining operation),
so any expression whose 23-op routing uses $k$ of these gates has a 16-op lower bound of
at least $k\cdot 3$ plus the non-EE/LL nodes. Verifying each case:
Form 1: 23-op = 2n; 16-op $\geq$ 5n (1 EEA = $3\times1=3n$ EEA-expansion + $1\ln=4n$ minimum). ✓
\end{proof}

\end{document}
